%====================================================================
% CS6007 Project Pitch (5 minutes)
% Federated Single-Index Bandits
%====================================================================
\documentclass[11pt]{beamer}
\usetheme{Madrid}
\usecolortheme{beaver}
\usefonttheme{professionalfonts}
\usepackage[T1]{fontenc}
\usepackage[utf8]{inputenc}
\usepackage{lmodern}
\usepackage{amsmath,amssymb,amsthm}
\usepackage{graphicx}
\usepackage{hyperref}
\usepackage{tikz}
\usetikzlibrary{arrows,positioning}
\usepackage{booktabs}
\usepackage{array}
\usepackage{pifont}

\setbeamertemplate{navigation symbols}{}
\setbeamerfont{footnote}{size=\tiny}
\newcommand{\yes}{\textcolor{green!45!black}{\ding{51}}}
\newcommand{\no}{\textcolor{red!70!black}{\ding{55}}}

\title[Multi-Agent Single-Index Bandits]{Multi-Agent Single-Index Bandits}
\subtitle{Learning unknown reward functions collaboratively under communication constraints}
\author[Harshitha \& Nischal]{\texorpdfstring{\textbf{Presented by:}\\[0.3em]\begin{tabular}{rl}\textcolor{blue}{Harshitha (23b1071)} & Dept. of CSE\\\textcolor{blue}{Nischal (23b1024)} & Dept. of CSE\end{tabular}\\[0.8em]\textbf{Course:} CS6007: Multi-Agent Machine Learning\\[0.2em]\textbf{Instructor:} Prof.\ Avishek Ghosh}{Harshitha, Nischal}}
\institute[IIT Bombay]{IIT Bombay}
\date{\small 28 September 2026}

\addtobeamertemplate{frametitle}{}{%
\begin{tikzpicture}[remember picture,overlay]
\node[anchor=north east,xshift=0.2cm,yshift=0.4cm] at (current page.north east){\includegraphics[height=1.1cm]{iitb_logo.png}};
\end{tikzpicture}}

\begin{document}

\usebackgroundtemplate{\tikz[remember picture,overlay]\node[opacity=0.15] at ([yshift=-2cm]current page.center){\includegraphics[width=0.6\paperwidth]{iitb_logo.png}};}
\begin{frame}\titlepage\end{frame}
\usebackgroundtemplate{}

%====================================================================
\section{Motivation}
%====================================================================
\begin{frame}{From contextual bandits to an unknown reward function}
\small
A contextual bandit chooses an action $x$ using its context and observes a noisy reward $y$.
\begin{block}{Common models}
\[
\mathbb{E}[y\mid x]=\langle x,\theta_*\rangle\quad\text{(Linear)},\qquad \mathbb{E}[y\mid x]=f(\langle x,\theta_*\rangle)\quad\text{(GLM)}.
\]
\end{block}
\begin{center}
\[
x\in\mathbb{R}^d\quad\longrightarrow\quad x^\top\theta_*\quad\longrightarrow\quad f(\cdot)\quad\longrightarrow\quad y
\]
\end{center}
\begin{alertblock}{The gap}
In GLMs, $f$ is known. In many applications, the reward curve is unknown and may be non-monotone.
\[
\boxed{\text{What if both }\theta_*\text{ and }f\text{ are unknown?}}
\]
\end{alertblock}
\end{frame}

%====================================================================
\begin{frame}{Single-Index Bandits: the problem}
\small
\begin{block}{Model}
At round $t$, the learner observes $K$ arms $\mathcal{X}_t=\{x_{t,a}\}_{a\in[K]}$, chooses $x_t$, and observes
\[
\boxed{y_t=f(\langle x_t,\theta_*\rangle)+\eta_t}
\]
where $\theta_*\in\mathbb{R}^d$ and $f:\mathbb{R}\rightarrow\mathbb{R}$ are both unknown, with $\|\theta_*\|_1=1$.
\end{block}
\begin{center}
\[
x_t\in\mathbb{R}^d\quad\xrightarrow{\ \theta_*\ }\quad x_t^\top\theta_*\quad\xrightarrow{\ f\ }\quad y_t
\]
\end{center}
\begin{alertblock}{Why is this hard?}
If $f$ is non-monotone,
\[
x_1^\top\theta_*>x_2^\top\theta_*\quad\not\Rightarrow\quad f(x_1^\top\theta_*)>f(x_2^\top\theta_*).
\]
\textbf{The largest projection need not be the best arm.}
\end{alertblock}
\end{frame}

%====================================================================
\begin{frame}{Key idea: estimate $\theta_*$ without knowing $f$}
\small
Define the score function
\[
S(x)=-\nabla\log p(x).
\]
Under Stein's identity,
\[
\mathbb{E}[yS(x)]=\mathbb{E}\!\left[f'(\langle x,\theta_*\rangle)\right]\theta_*.
\]
Let
\[
\mu_*=\mathbb{E}\!\left[f'(\langle x,\theta_*\rangle)\right].
\]
Then
\[
\boxed{\mathbb{E}[yS(x)]=\mu_*\theta_*}
\]
so the direction of
\[
\boxed{\hat v=\frac{1}{T_0}\sum_{t=1}^{T_0}y_tS(x_t)}
\]
estimates the direction of $\theta_*$ without estimating $f$.
\begin{block}{Key insight}
\centering
\textbf{Unknown $f$ does not prevent us from estimating the index direction.}
\end{block}
\end{frame}

%====================================================================
\begin{frame}{The SIB problem reduces to 1-D learning}
\small
Once $\theta_*$ is estimated, project each context onto the estimated direction:
\[
z=x^\top\hat\theta.
\]
\begin{center}
\[
x\in\mathbb{R}^d\quad\xrightarrow{\ \hat\theta\ }\quad z=x^\top\hat\theta\in\mathbb{R}\quad\xrightarrow{\ f\ }\quad f(z)
\]
\end{center}
The original $d$-dimensional problem becomes learning an unknown function on a 1-D index.
\[
[-1,1]=I_1\cup I_2\cup\cdots\cup I_M,\qquad M=O(T^{1/3}).
\]
\begin{block}{ZoomSIB-UCB}
The index domain is divided into bins, and a sleeping-UCB procedure learns which regions of the unknown reward curve are promising.
\end{block}
\end{frame}

%====================================================================
\section{Federation}
%====================================================================
\begin{frame}{Problem Formulation: Multi-Agent SIB}
\small
\begin{block}{Setting}
Consider $N$ decentralized agents on a connected communication graph $G$. At each round $t$, each agent $i$ observes an independent arm set, pulls $x_{i,t}$, and observes $y_{i,t}=f(\langle x_{i,t},\theta_*\rangle)+\eta_{i,t}$. The parameter $\theta_*$ and link $f$ are identical across agents but unknown. Agents can exchange messages with their neighbors subject to a communication budget.
\end{block}
\begin{block}{Objective: Cumulative Network Regret}
\[
R_T(N)=\sum_{t=1}^{T}\sum_{i=1}^{N}\left[f(\langle x^*_{i,t},\theta_*\rangle)-f(\langle x_{i,t},\theta_*\rangle)\right].
\]
\end{block}
\begin{center}
\emph{Can decentralized collaboration improve per-agent regret when neither the index direction nor the link function is known?}
\end{center}
\end{frame}

%====================================================================
\begin{frame}{Fed-ZoomSIB: federate both phases}
\small
\begin{columns}[T]
\begin{column}{0.48\textwidth}
\begin{block}{Phase 1: Estimate $\theta_*$}
Each agent computes
\[
\hat\theta_i\propto\frac{1}{T_0}\sum_t y_{i,t}S(x_{i,t}).
\]
Then agents exchange their local estimates:
\[
\boxed{\hat\theta=\frac{1}{N}\sum_{i=1}^{N}\hat\theta_i}
\]
\textbf{Exact pooled averaging}\\[-0.2em]
One $d$-dimensional message.
\end{block}
\end{column}
\begin{column}{0.48\textwidth}
\begin{block}{Phase 2: Learn $f$}
Project:
\[
z_{i,t}=x_{i,t}^{\top}\hat\theta.
\]
Each agent maintains bin statistics
\[
(n_j,S_j).
\]
Agents periodically exchange these statistics.
\textbf{Statistical pooling.}
\end{block}
\end{column}
\end{columns}
\vspace{0.2em}
\begin{center}
\boxed{\text{Phase 1: exact pooling}\qquad\text{vs.}\qquad\text{Phase 2: statistical pooling}}
\end{center}
\end{frame}

%====================================================================
\begin{frame}{What do we want to find out?}
\small
\begin{block}{1. Does collaboration help?}
Compare per-agent regret under independent and federated learning:
\[
R_T^{\mathrm{federated}}\quad\text{vs.}\quad R_T^{\mathrm{independent}}.
\]
\end{block}
\begin{block}{2. How much communication is enough?}
Study the tradeoff
\[
\boxed{\text{Regret}\quad\leftrightarrow\quad\text{Communication}}
\]
using periodic, event-triggered, and quantized exchanges.
\end{block}
\begin{block}{3. How does the network matter?}
Study dependence on the number of agents, dimension, number of arms, and graph connectivity:
\[
N,\qquad d,\qquad K,\qquad 1-|\lambda_2(\mathbf P)|.
\]
\end{block}
\end{frame}

%====================================================================
\section{Evaluation}
%====================================================================
\begin{frame}{Evaluation Plan}
\scriptsize
\renewcommand{\arraystretch}{1.1}
\begin{tabular}{>{\raggedright\arraybackslash}p{1.8cm}>{\raggedright\arraybackslash}p{5.2cm}>{\raggedright\arraybackslash}p{3.0cm}}
\toprule
\textbf{Dates}&\textbf{Work}&\textbf{Deliverable}\\
\midrule
29 Sep -- 10 Oct&Reproduce \textsc{ZoomSIB-UCB} and baselines: Random, LinUCB, ESTOR, GSTOR, IGP-UCB&Validated regret curves\\
\addlinespace
11 -- 22 Oct&Implement Fed-ZoomSIB: Stein averaging + cooperative sleeping-UCB&Federated speedup curves\\
\addlinespace
23 -- 29 Oct&Communication ablations: periodic sync, event-triggered updates, quantization&Regret vs.\ communication\\
\addlinespace
30 Oct -- 5 Nov&One extension: safety or clustered index directions&Extension analysis\\
\addlinespace
5 -- 12 Nov&Benchmarking, final report, and presentation&Project submission\\
\bottomrule
\end{tabular}
\vspace{0.3em}
\textbf{Metrics:}\quad Network regret \quad Per-agent regret \quad Messages / bytes
\[
\boxed{\text{Vary }N,\ d,\ K,\ \text{topology, and link shape}}
\]
\end{frame}

%====================================================================
\begin{frame}{Where does our problem sit?}
\centering
\renewcommand{\arraystretch}{1.35}
\begin{tabular}{>{\centering\arraybackslash}p{3cm}>{\centering\arraybackslash}p{3cm}>{\centering\arraybackslash}p{3cm}}
\toprule
&\textbf{Single Agent}&\textbf{Multi-Agent}\\
\midrule
\textbf{Known reward structure}&Linear / GLM Bandits&Cooperative Contextual Bandits\\
\midrule
\textbf{Unknown reward function}&SIB / ZoomSIB-UCB&\alert{\textbf{Our problem}}\\
\bottomrule
\end{tabular}
\vspace{0.6em}
\[
\boxed{\text{Unknown link}+\text{Multi-agent cooperation}=\text{Federated SIB}}
\]
\vspace{0.2em}
\small
Our focus is the interaction between statistical learning and communication constraints.
\end{frame}

%====================================================================
\begin{frame}{References}
\footnotesize
\begin{enumerate}
\item D.~Dey, S.~Bhore, A.~Ghosh. \emph{Optimal Regret for Single Index Bandits.} arXiv:2605.09454, 2026.
\item Y.~Kang, M.~Liu, B.~Yi, J.~Lyu, Z.~Zhang, D.~Zhou, Y.~Li. \emph{Single Index Bandits: Generalized Linear Contextual Bandits with Unknown Reward Functions.} ICLR, 2026.
\item A.~Dubey, A.~Pentland. \emph{Kernel Methods for Cooperative Multi-Agent Contextual Bandits.} ICML, 2020.
\item S.~Amani, C.~Thrampoulidis. \emph{Decentralized Multi-Agent Linear Bandits with Safety Constraints.} AAAI, 2021.
\item S.~Arya, S.~Bhattacharjee, B.~K.~Sriperumbudur. \emph{Kernel Single-Index Bandits: Estimation, Inference, and Learning.} arXiv:2603.18938, 2026.
\item A.~Ghosh, J.~Chung, D.~Yin, K.~Ramchandran. \emph{An Efficient Framework for Clustered Federated Learning.} NeurIPS, 2020.
\item A.~Ghosh, S.~R.~Chowdhury, A.~Gopalan. \emph{Misspecified Linear Bandits.} AAAI, 2017.
\end{enumerate}
\end{frame}

%====================================================================
% thankyou.tex
% thankyou.tex

\begin{frame}[plain]
  % light colored background
  \setbeamercolor{background canvas}{bg=blue!10}
  \centering
  \vfill
  {\Huge\bfseries\color{blue!80!black} Thank You!}
  \vfill
\end{frame}
%====================================================================
\end{document}